% Mục 4.5: Mutation testing, trả lời RQ1. Nguồn: results/{apple-m5,C4M7Y670NT,LAP1184}/mutation.log (toàn bộ).
% Số episode lúc bắt lỗi: results/apple-m5/mutation_logs/barrier_mut2.run{2,5}.log,
% barrier_mut2_fuzz.run{1..5}.log; barrier_mut4.run1.log, barrier_mut4_fuzz.run1.log;
% results/C4M7Y670NT/mutation_logs/barrier_mut4.run{1,2,4,5}.log.

\label{sec:corr:mutation}
\begin{rqanswer}\textbf{Trả lời RQ1.} Không bài kiểm tra nào tìm thấy lỗi trên ba máy, và bộ kiểm tra bắt
cả 14 biến thể hỏng. Tính đúng dựa trên lập luận ở mục Cài đặt. Kiểm thử củng cố lập luận đó. Trên
\PCname{apple-m5}, \PCname{C4M7Y670NT} và \PCname{LAP1184}, mọi cài đặt đạt mọi bài kiểm tra (12 lock,
7 barrier và 19 counter trên hai máy Apple, 13, 9 và 20 trên \PCname{LAP1184}), và TSan không báo data
race nào. Hai lỗi memory order khó bắt nhất. Stress test không bắt \texttt{TAS\_RELEASE\_RELAXED} lần nào
(0/3 trên cả ba máy). Bản thường bắt \texttt{fetch\_add\_release\_only} 1/5 lần trên \PCname{apple-m5},
4/5 trên \PCname{C4M7Y670NT} và 0/5 trên \PCname{LAP1184}. TSan bắt cả hai lỗi trên mọi máy.\end{rqanswer}
% nguồn: results/{apple-m5,C4M7Y670NT,LAP1184}/mutation.log: TAS_RELEASE_RELAXED MISSED 0/3, TSan CAUGHT (cả ba);
% fetch_add_release_only plain caught 1/5 (apple-m5), 4/5 (C4M7Y670NT), SURVIVED 0/5 (LAP1184), DATA RACE reported (cả ba);
% results/{máy}/test.log, tsan.log (số cài đặt, 0 cảnh báo).

\emph{Bằng chứng.} Bảng~\ref{tab:corr:mutation} tóm tắt kết quả trên \PCname{apple-m5}. Biến thể lock là
macro sửa một chỗ trong \texttt{lock.h}. Biến thể barrier và counter là bản sao trong \texttt{tests/},
mỗi bản cài một lỗi. Riêng \texttt{single\_flag} và \texttt{block\_reserve} là thuật toán khác hẳn. Khi
chạy biến thể hỏng, watchdog báo treo sau 5~s với lock và 2~s với barrier, và mỗi pha barrier có ngân
sách 0,5~s.
% nguồn: include/lock.h (PC_MUTANT_TAS_*, PC_MUTANT_MCS_*), tests/test_barrier.cpp (namespace mutant, PC_MUTANT 1..5),
% tests/test_counter.cpp (PC_MUTANT_plain_long..block_reserve); mk/part1.mk:44 (watchdog=5);
% mk/part2.mk:63 (stall_s=2 cap_s=0.5).

\begin{table}[tbp]\centering\scriptsize\setlength{\tabcolsep}{3pt}
% nguồn: results/apple-m5/mutation.log
\begin{tabularx}{\linewidth}{@{}>{\ttfamily}l>{\raggedright\arraybackslash}Xlll@{}}
\toprule
\normalfont\textbf{Biến thể} & \textbf{Lỗi cài vào} & \textbf{Stress (+ fuzzing)} & \textbf{TSan} & \textbf{Bị bắt bởi} \\
\midrule
TAS\_CHECK\_THEN\_SET & \texttt{load} rồi \texttt{store} thay \texttt{exchange} & 3/3 vi phạm loại trừ & race & cả hai \\
TAS\_RELEASE\_RELAXED & nhả lock bằng \texttt{store} relaxed & 0/3 & race & chỉ TSan \\
MCS\_NO\_WAIT\_SUCCESSOR & nhả lock không chờ thread kế tiếp & 3/3 treo & không & stress \\
MCS\_ALWAYS\_CLEAR\_TAIL & luôn đặt \texttt{tail} về null & 3/3 treo & race & cả hai \\
\midrule
single\_flag & một cờ, không đảo chiều & 5/5 (5/5) sai episode & không & stress \\
flip\_before\_reset & đảo \texttt{sense} trước khi đặt lại bộ đếm & 2/5 (5/5) treo & không & stress \\
nonatomic\_count & \texttt{load}+\texttt{store} thay \texttt{fetch\_add} & 5/5 (5/5) treo & không & stress \\
fetch\_add\_release\_only & \texttt{fetch\_add} release thay acq\_rel & 1/5 (0/5) sai episode & race & chủ yếu TSan \\
thread\_local\_sense & local sense chung cho mọi barrier & 5/5 (5/5) sai giá trị handoff & race & cả hai \\
\midrule
plain\_long & \texttt{long} thường & 3/3 final, dup, miss, order & race & cả hai \\
load\_store & \texttt{load} rồi \texttt{store} atomic & 3/3 final, dup, miss, order & không & kiểm giá trị \\
cas\_no\_retry & CAS một lần, không thử lại & 3/3 final, dup, miss & không & kiểm giá trị \\
lock\_read\_outside & đọc giá trị trước \texttt{acquire()} & 3/3 final, dup, miss, order & race & cả hai \\
block\_reserve & giữ trước 8 giá trị mỗi thread & 3/3 chỉ real-time & không & kiểm giá trị \\
\bottomrule
\end{tabularx}
\caption{Mutation testing trên \PCname{apple-m5}: lock (trên), barrier (giữa), counter (dưới). Cột stress
cho số lần chạy bắt được trên số lần chạy (trong ngoặc: bản schedule fuzzing) cùng dấu hiệu hoặc phép kiểm
báo lỗi; cột TSan cho biết TSan có báo data race không.}
\label{tab:corr:mutation}
\end{table}

Trên \PCname{C4M7Y670NT}, mọi biến thể bị bắt bởi cùng kỹ thuật như trên \PCname{apple-m5}. Chỉ số lần
stress test bắt được khác. Ở bản thường, \texttt{fetch\_add\_release\_only} bị bắt 4/5 lần (so với
1/5) và \texttt{flip\_before\_reset} 1/5 lần (so với 2/5). Bản fuzzing cho cùng kết quả như trên
\PCname{apple-m5} (0/5 và 5/5). Kết quả TSan của máy này chỉ dựa trên dòng tóm tắt trong
\texttt{mutation.log}, vì log chi tiết của các lần chạy TSan không được lưu.
% nguồn: results/C4M7Y670NT/mutation.log (Part 2: flip_before_reset caught 1/5 | caught 5/5;
% fetch_add_release_only caught 4/5 | SURVIVED 0/5 | DATA RACE reported); các dòng còn lại trùng apple-m5.
% results/C4M7Y670NT/mutation_logs/ không có *_tsan.log, counter_muttsan_*, locks_muttsan_*.

Trên \PCname{LAP1184}, mọi biến thể bị bắt bởi cùng kỹ thuật như trên \PCname{apple-m5}, với ba khác
biệt. Thứ nhất, \texttt{fetch\_add\_release\_only} không bị stress test bắt lần nào (0/5, fuzzing 0/5),
chỉ TSan bắt. Thứ hai, bản thường bắt \texttt{flip\_before\_reset} 5/5 lần, so với 2/5 trên
\PCname{apple-m5}. Thứ ba, với \texttt{thread\_local\_sense}, bản TSan dừng ngay khi phép kiểm handoff
báo sai giá trị, và TSan không in cảnh báo nào trước đó. Trên máy này biến thể chỉ bị phép kiểm giá trị
bắt (5/5 ở bản thường và bản fuzzing).
% nguồn: results/LAP1184/mutation.log (Part 2: flip_before_reset caught 5/5 hang 5; fetch_add_release_only
% SURVIVED 0/5 | SURVIVED 0/5 | DATA RACE reported; thread_local_sense: no report; checks: wrong
% episode/value; Part 1: TAS_RELEASE_RELAXED MISSED 0/3, TSan CAUGHT);
% results/LAP1184/mutation_logs/barrier_mut5_tsan.log (FAIL [handoff]: WRONG VALUE, không có dòng TSan).

\emph{Vì sao.} Hai biến thể chỉ TSan bắt ổn định đều là lỗi memory order. \texttt{TAS\_RELEASE\_RELAXED}
đổi lệnh nhả lock \texttt{stlrb} (store release) thành \texttt{strb}. TSan bắt nó vì store relaxed không
tạo cạnh happens-before giữa hai chủ lock liên tiếp. Ba lần chạy stress không có cập nhật nào bị mất.
Nhóm chưa đo vì sao lỗi không lộ trên ARM. Một khả năng là các phép ghi trong critical section thường
đã tới cache trước khi phép ghi nhả lock được thấy, nên cửa sổ lỗi rất hẹp.
\texttt{fetch\_add\_release\_only} đổi \texttt{ldaddal} thành \texttt{ldaddl}, nên thread tới cuối có
thể đọc ô của thread khác trước khi \texttt{fetch\_add} của nó có hiệu lực. TSan bắt nó qua các ô thường
của phép kiểm handoff.

Trên x86-64, mọi \texttt{fetch\_add} đều là lệnh \texttt{lock xadd}. Mọi store, kể cả \texttt{relaxed}
và \texttt{release}, đều là \texttt{mov}. Vì vậy trên \PCname{LAP1184} hai biến thể memory order chạy đúng
lệnh máy của bản đúng, và không lịch chạy nào làm lộ lỗi.
% nguồn: biên dịch include/lock.h, tests/test_barrier.cpp (bản đúng và hai biến thể) bằng Apple clang
% -arch x86_64 -O2 -S: assembly của biến thể trùng bản đúng. LAP1184 dùng g++ 11.5, chưa so assembly riêng.

Ngược lại, TSan im lặng với bảy biến thể chỉ dùng atomic hoặc gây treo. Không kỹ thuật nào một mình
bắt được mọi biến thể. Stress test hoặc phép kiểm giá trị bắt mọi lỗi thuật toán. Bốn lỗi treo (hai biến
thể MCS, \texttt{flip\_before\_reset}, \texttt{nonatomic\_count}) hiện ra qua watchdog.

Fuzzing làm \texttt{flip\_before\_reset} treo sau 3--7 episode (bản thường cần 39\,896 và 289\,052
episode ở hai lần bắt được trên \PCname{apple-m5}). Bản fuzzing không bắt được lỗi thiếu acquire (0/5).
Bản này chỉ chạy một phần mười số episode, và trên hai máy Apple bản thường chỉ bắt được lỗi sau hơn
13\,000 episode của một pha. Nhóm chưa đo vì sao bản thường trên \PCname{LAP1184} bắt
\texttt{flip\_before\_reset} thường hơn. Một khả năng là máy ảo tước core của thread nhiều hơn.
% nguồn: results/apple-m5/mutation_logs/barrier_mut2.run{2,5}.log (39 896, 289 052),
% barrier_mut2_fuzz.run{1..5}.log (5, 7, 7, 3, 5 episode); mutation.log (fetch_add_release_only fuzzing 0/5).
% Episode lúc bắt fetch_add_release_only ở bản thường: apple-m5 barrier_mut4.run1.log 35 881 (pha reuse);
% C4M7Y670NT barrier_mut4.run{1,2,4,5}.log 31 763, 13 373, 23 689, 26 190 (pha all-cores).
% Bản fuzzing: g_size = 0.1 (tests/test_barrier.cpp, main), pha reuse 30 000 và all-cores 10 000 episode
% (apple-m5/mutation_logs/barrier_mut4_fuzz.run1.log).

\emph{Độ vững.} Kết luận trên dựa trên một lượt chạy mỗi máy và chỉ trên các lịch chạy đã xảy ra.
Mục~\ref{sec:corr:limits} nêu các giới hạn này.
